\section{Discussion}\label{sec:discussion}

Table~\ref{tab:indicators-summary} sets the four tests of Sections~\ref{sec:results} and~\ref{sec:phase1} side by side. Each row asks the same question: does the test tell a learner that has learned a continuation from one that carries a stale estimate, or does it merely record that the memory-1 learner can condition on rivals' bids and the memory-0 learner cannot?

\begin{table}[t]
\caption{What each indicator detects, and whether it separates the two learners. Notes: memory-1 column is the long (5M-round) run, 100 seeds; memory-0 column is the long memory-0 control, 100 seeds; ``discriminates'' asks whether the test tells learned continuation from a stale estimate, not merely whether the two columns differ.}
\label{tab:indicators-summary}
\centering
\footnotesize
\setlength{\tabcolsep}{4pt}
\begin{tabular}{p{1.6cm} p{3.2cm} p{3.6cm} p{3.0cm} p{2.9cm}}
\toprule
Test & Meant to detect & Memory-1 result & Memory-0 control & Discriminates? \\
\midrule
M3 punishment & Rivals answer a deviation, then forgive & Rivals react in 91\% of seeds (deviator firm 1); punish-forgive 21\%, no recovery 70\%, no reaction 9\% & No reaction in 100\% of seeds & No: it detects memory, which a one-state table lacks by construction \\
M4 ablation & Price needs memory and foresight & $\Dl = 0.621$; $\gamma = 0$ gives $0.370$, memory 0 gives $0.943$; neither below the $0.25$ threshold & Is the ablation & No \\
M5 temptation & Profitable deviation left unplayed & 100\% of seeds; mean gap $\eur{2{,}226}$ per round & 100\% of seeds; mean gap $\eur{1{,}448}$ per round & No: fires for the memory-0 control \\
Bellman residual & One fresh update would flip the decision & 31\% of temptations; 12 seeds at 0\%, none at 100\% & 100\% of temptations in every seed & Yes: 0.31 versus 1.00, no overlap across seeds \\
\bottomrule
\end{tabular}
\end{table}

\paragraph{A reading of the memory-0 price.}
What follows is a verbal argument, not a theorem; the evidence for it is the occupation measure of Section~\ref{sec:occupation} and the robustness results of Sections~\ref{sec:beta}, \ref{sec:reinject} and~\ref{sec:payasbid}. Because the clearing price is the second-lowest bid (Section~\ref{sec:market}), a lone undercut from a profile whose two lowest bids are tied sells 60 MW at the rivals' price and leaves the price where it was; at the all-cap profile that is a gain of $\eur{2{,}400}$ per round. This is why the M5 gap is large at every supra-competitive profile, and why the deviation round in M3 is profitable for the deviator; whether the deviation pays over the following rounds depends on the rivals' reaction and is not what M3 measures. The price falls only when a second firm follows the first down. On this reading a race to the bottom needs two firms to keep undercutting each other, and it stops wherever it stands when the exploratory draws that drive it become rare; the observed rest point, one firm at or near cost and two tied near the cap, is what such a stalled race looks like under this clearing rule. The walk down is a chain of learning events, each requiring a specific exploratory draw at a specific time; the move up from the bottom, where a firm has nothing to lose by raising its bid, needs a single event. Figure~\ref{fig:mechanism} draws the asymmetry. For the memory-0 learner the process spends its time near the top of the grid, more so the less exploration there is, and freezing the tables samples that distribution. The occupation measure in Section~\ref{sec:occupation} supports the last step of this argument, though not the asymmetry itself: the memory-0 time average rises from $\Dl = 0.446$ at $\varepsilon = 0.1$ to $0.934$ at $\varepsilon = 0.001$, within $0.01$ of the frozen $0.943$. The analogy is stochastic stability \citep{kandori1993,young1993}: as noise vanishes such a process selects the states that take the most rare draws to leave and the fewest to reach. We have computed no escape probabilities, so this is an analogy rather than a result, but the high-price profiles are where the memory-0 process is found. On this reading the memory-0 process is an Edgeworth-type cycle run by a learning rule \citep{maskin1988} rather than an instance of the folk theorem \citep{fudenberg1986}: with one state there is no memory to carry a threat. The memory-1 learner's frozen value is also the small-noise limit of its occupation measure, but a far more fragile one: its exploring process sits below that value until $\varepsilon = 10^{-4}$, roughly 30 times less noise than memory-0 needs, which is consistent with a coordinated cycle that single random draws break. The occupation measure separates the two learners by noise tolerance, not by kind (Section~\ref{sec:occupation}).

\begin{figure}[t]
\centering
\begin{tikzpicture}[>=Stealth, font=\footnotesize, x=1cm, y=1cm]
% the bid ladder: six rungs stand for the fifteen
\foreach \i in {0,...,5} {\draw[gray!70] (0,\i*1.1) rectangle ++(1.1,1.1);}
\node[above=1pt] at (0.55,6.6) {$p_{15}=\eur{100}$ (cap)};
\node[below=1pt] at (0.55,0) {$p_{1}=\eur{10}$ (cost)};
\node[gray] at (0.55,3.85) {$\vdots$};
% left: the walk down, one completed chain per rung
\foreach \i in {5,4,3} {\draw[->, thick] (-0.35,\i*1.1+0.55) -- (-0.35,\i*1.1-0.55);}
\draw[->, thick, dotted] (-0.35,2.75) -- (-0.35,0.55);
\draw[gray] (-0.75,6.05) -- (-0.75,3.35);
\draw[gray] (-0.75,6.05) -- (-0.6,6.05);  \draw[gray] (-0.75,3.35) -- (-0.6,3.35);
\node[anchor=east, text width=4.6cm, align=right] at (-0.9,4.7)
  {\textbf{Walk down: one completed chain per rung.}};
\node[anchor=east, text width=4.6cm, align=right, gray] at (-0.9,1.6)
  {\ldots\ and again for every rung, all the way down.};
% right: the jump up, one event
\draw[->, very thick] (1.45,0.55) -- (1.45,6.05);
\node[anchor=west, text width=3.2cm, align=left] at (1.7,3.3)
  {\textbf{Jump up: one event.}\\ A firm bidding at cost earns nothing on what it sells, so when the price is at the bottom raising its bid costs it nothing and needs no chain.};
% the consequence
\node[draw, rounded corners, text width=4.2cm, align=left, inner sep=5pt, anchor=west] at (5.4,5.2)
  {\textbf{Where the process lives:} near the top, more so the less exploration there is.};
\end{tikzpicture}
\caption{A verbal reading of the memory-0 price (Section~\ref{sec:discussion}); it is not a theorem, and the asymmetry it draws was not measured directly. Rungs are the bid grid. On the left, each step down the grid would need a completed chain of learning events, each link an exploratory draw at the right moment. On the right, a firm at the bottom would raise its bid in one event, having nothing to lose. If the asymmetry holds it would concentrate the exploring process near the top; what Section~\ref{sec:occupation} does measure is that the process is found there, and that freezing samples that distribution.}
\label{fig:mechanism}
\end{figure}

\paragraph{What the indicators can and cannot do.}
The three behavioural tests were designed for a learner whose high price rests on anticipated retaliation. Here they cannot tell such a learner from one whose price is a property of the learning rule. M5 fires in 100\% of seeds for both learners: the memory-0 control leaves the same kind of money on the table, so the fingerprint does not require a learner that can anticipate anything. M4 fails because removing memory raises the price to $0.943$ rather than lowering it, and removing foresight leaves $0.370$, above the collapse threshold of $0.25$. M3 does separate the two, but for a trivial reason: a one-state table has nothing to react to, so it records no reaction in 100\% of memory-0 seeds. It confirms that the memory-1 learner's greedy policy varies with rivals' last bids, which its state definition permits but does not guarantee, and its verdict for a given seed changes with the identity of the deviator in 61 of 100 seeds. The residual is different in kind. It asks whether one fresh update of the deviation's value would flip the decision. For the memory-0 learner the answer is yes at every temptation in every seed. For the memory-1 learner it is yes at 31\% of temptations, no seed reaches 100\%, and twelve seeds sit at 0\%. In the remaining 69\% of memory-1 temptations the continuation after deviating, valued by the learner's own table, is low enough that a single fresh backup would leave the decision unchanged. That is a learned continuation, the thing tacit collusion consists of, measured on the table rather than inferred from behaviour. Whether a learned continuation of this kind deserves the name is a question this paper does not settle, and the measurement has a limit of its own: the continuation value it uses comes from the same table and may itself be stale (Section~\ref{sec:residual}).

Two consequences follow for anyone screening simulated bidding algorithms. First, in this market a high $\Dl$ together with a large unclaimed temptation does not by itself distinguish a learner that anticipates retaliation from one that cannot: the memory-0 control produces both. Whether the same holds in other non-pivotal auctions was not tested. Second, the residual needs the Q-table. It is a simulation-side diagnostic, available to a researcher, or to a firm auditing its own algorithm if that algorithm keeps a value table. None of this is a screen for real bid data, where only bids, prices and quantities are observed and no table exists to be read.

\paragraph{Relation to the literature.}
\citet{abada2023} and \citet{abada2024} reach a similar broad conclusion in a different market: outcomes that look collusive can arise from how independent learners explore. This paper agrees for the memory-0 learner and adds that the memory-1 learner is only partly of that kind. The mechanism detail here is different. Their concern is that learners fail to find the competitive outcome; here the freeze is not specific to the uniform-price rule, since it survives the change to pay-as-bid (Section~\ref{sec:payasbid}), though under uniform pricing the rule shapes where the process stops: at the (low, high, high) rest point the low bidder is paid the high price for its full 60 MW and has no profitable deviation, while each high bidder could earn \eur{1{,}543} more per round by shaving one rung (Section~\ref{sec:market}) but can discover that only through an exploratory draw. \citet{asker2022} find that the learning protocol, in particular what the algorithm is told about profits it did not earn, determines whether prices end high or low in their game. The memory-0/memory-1 contrast is a difference of state space rather than of updating rule, so it is not the protocol difference they study; what the two papers share is that a design choice inside the learner, not the incentives of the game, moves the outcome. The residual share separated the two learners' tables in this market; whether it would separate other protocol differences was not tested. \citet{banchio2022} find that the auction format decides whether Q-learning bidders settle at collusive-looking bids: below value under first-price rules, not under second-price rules. In this market the format moved the rest point but did not remove the freeze: under pay-as-bid the memory-0 $\Dl$ falls from $0.943$ to $0.845$, the most common frozen profiles change from one low bidder with two tied high bidders to symmetric ties, and every seed still converges with $\Dl$ above $0.5$.

\paragraph{Limits.}
The study is a simulation of one stylised single-node market with symmetric firms, inelastic and constant demand, no network or congestion, no multi-segment bid curves, and tabular learners only, so nothing is said about function-approximation learners. Apart from $\beta$ and the horizon the hyperparameters were not swept; $\alpha$, $\gamma$ and the grid come from the pre-registered specification. The punishment verdict uses fixed 12-round windows and a half-rung threshold; the nine long-run seeds with 5- or 7-cycles can be misread by those windows. The Bellman residual is a one-step backup with the current table, so $\max_{a'} Q_i(s', a')$ may itself be stale; a $k$-step rollout under the greedy policy would be stricter and was not done. The memory-1 occupation curve has one point at $\varepsilon = 10^{-4}$, so the transition is not located precisely. There is no theory: the Edgeworth-cycle argument is verbal. Finally, the two extra one-shot Nash equilibria on the grid ($\Dl = 0.07$ and $0.14$) mean the competitive benchmark is a band $[0, 0.14]$ rather than a point; all reported $\Dl$ lie far above it.

\paragraph{Next steps.}
A $k$-step rollout residual would remove the dependence on possibly stale continuation values; whether the separation survives it is an open question. A run at an exploration rate between $\varepsilon = 10^{-3}$ and $10^{-4}$ would locate where the memory-1 occupation measure meets its frozen value. A two-firm, three-price version of the market might be solvable by hand and might turn the verbal cycle argument into a statement about basins of the noisy dynamic. Network constraints and multi-segment bids are separate future work: they change the pivotality on which the argument above rests and should be studied on their own.
